\section{Summary of Outcomes}
\label{sec:summary_outcomes}

This chapter consolidates the empirical outcomes of the \texttt{DetectorNX-G3-SNN} evaluation against the project's success criteria, validating the neuromorphic paradigm for safety-critical edge deployment.

\subsection{Global Precision Parity}
The primary outcome of the spatial evaluation was the achievement of a \gls{miou} of $0.4613$ for the spiking model, representing $94.9\%$ of the \gls{cnn} baseline ($0.4861$), and an \gls{map_50} of $0.5977$ against the \gls{cnn}'s $0.6550$ --- a $91.3\%$ parity under the stricter COCO-convention detection-quality metric.
Furthermore, both architectures achieved a perfect recall of $100.0\%$ for objects with an \gls{iou} $> 0.5$.
These results prove that the architectural refinements introduced in this work, specifically the multi-box architecture, \gls{se} and \gls{sew} residual blocks, successfully mitigated the precision loss inherent in $1$-bit quantization and discrete temporal resolution.

\subsection{Decisive Efficiency Gains}
The transition from dense \gls{cnn} convolutions to event-driven spiking dynamics yielded an average power saving of $94.33\%$ across the validation subset,
which was driven by a $3.45\times$ reduction in the total operations required to process a high-definition frame ($9.22$ G-FLOPs vs. $2.67$ G-SOPs). Crucially, the \gls{snn} demonstrated ``Reactive Scalability'' ($R^2=0.7514$), where power consumption scaled linearly with scene complexity rather than model size,
a property that positions neuromorphic perception as inherently
suited to the intermittent, sparse event streams characteristic
of real-world automotive environments.
Benchmarked against the contemporary published landscape of \gls{snn} object detectors (Table~\ref{tab:efficiency_landscape}), \texttt{DetectorNX-G3-SNN} achieves the highest \gls{map_50}/mJ ratio in the surveyed cohort, leading the next-best entry by approximately $8\times$. The dataset asymmetry caveat applies, but the magnitude of the lead is robust to conservative multi-class adjustment.

\subsection{Temporal Reliability and Sharpening}
Beyond global averages, the spiking paradigm demonstrated a unique ``Bounding Box Sharpening'' effect. As temporal evidence accumulated across $T=10$ timesteps, the \gls{snn} was observed to produce more stable and spatially accurate regressions in high-traffic scenes compared to the single-pass \gls{cnn}. 
This temporal integration capability allowed the \gls{snn} to exhibit a more decisive confidence distribution, being confident enough to classify noise with a confidence of $0\%$, effectively suppressing low-evidence background noise.

\subsection{Real-Time Viability}
Despite the hardware-mismatch penalty of executing iterative spiking logic on synchronous \gls{gpu} hardware, the \gls{snn} achieved a mean inference latency of $3.85$\,ms ($259.80$ \gls{fps}). 
When evaluated against the industry-standard threshold for real-time automotive perception of $30$ \gls{fps} \cite{khan2024real}, the spiking model maintains an $8.6\times$ safety margin, despite using simulated \gls{lif} nodes. 
This validates the immediate viability of the proposed architecture for live deployment in vehicle-to-everything (V2X) and driver assistance systems.

\subsection{Success Criteria Validation}
\label{subsec:sc_validation}

Table~\ref{tab:sc_outcomes} consolidates the empirical outcome of each success criterion defined in Section~\ref{sec:intro_success_criteria} (mapped to its validating experiment in Table~\ref{tab:sc_traceability}), confirming that all six criteria were satisfied.

\begin{table}[htbp]
\centering
\resizebox{\textwidth}{!}{%
\begin{tabular}{|c|p{4.5cm}|p{6cm}|c|}
\hline
\textbf{SC} & \textbf{Quantitative Target} & \textbf{Achieved Result} & \textbf{Status} \\
\hline
SC1 & \gls{miou} and \gls{map_50} $\geq 90\%$ of \gls{cnn} baseline & $94.9\%$ \gls{miou} parity, $91.3\%$ \gls{map_50} parity & Satisfied \\
\hline
SC2 & Reduction in theoretical energy per inference; reduction realised on \gls{gpu} hardware & $94.33\%$ mean theoretical saving on the validation subset; $65.12\%$ extrapolated on Intel Loihi & Satisfied \\
\hline
SC3 & $\geq 30$ \gls{fps} on \gls{gpu} hardware & $259.80$ \gls{fps} ($3.85$\,ms latency); $8.6\times$ over threshold & Satisfied \\
\hline
SC4 & Parameter parity within $1\%$ & \gls{snn} $11.72$M vs \gls{cnn} $11.73$M ($0.09\%$ delta) & Satisfied \\
\hline
SC5 & $100\%$ recall at \gls{iou} $> 0.5$ & $100.0\%$ recall (both architectures) & Satisfied \\
\hline
SC6 & \gls{snn} confidence calibration $\geq$ \gls{cnn} & \gls{snn} superior at low confidence (suppresses false positives); qualitatively equivalent at high confidence & Satisfied \\
\hline
\end{tabular}%
}
\caption[Success criteria validation outcomes.]{Empirical outcomes for each success criterion defined in Section~\ref{sec:intro_success_criteria}. Validating experiments are mapped in Table~\ref{tab:sc_traceability}.}
\label{tab:sc_outcomes}
\end{table}

\vspace{1em}
In summary, the empirical evidence demonstrates that the \texttt{DetectorNX-G3-SNN} has successfully bridged the precision gap, establishing the spiking paradigm not merely as a theoretical curiosity, but as a robust, first-class solution for high-definition, power-constrained automotive perception.
